\documentclass[10pt, a4paper]{article}
\usepackage[utf8]{inputenc}
\usepackage[norsk]{babel}
\usepackage{amsmath, amssymb}
\usepackage{geometry}
\usepackage{xcolor}
\usepackage{enumitem}
\usepackage[most]{tcolorbox}
\usepackage{multicol}
\usepackage{csquotes}
\usepackage{eurosym}

\geometry{margin=1.5cm, top=1.2cm, bottom=1.2cm}

\definecolor{s1color}{RGB}{0,102,204}
\definecolor{s2color}{RGB}{0,153,76}
\definecolor{s3color}{RGB}{204,0,0}
\definecolor{boxbg}{RGB}{245,245,250}
\definecolor{boxborder}{RGB}{80,80,120}
\definecolor{formulabg}{RGB}{255,250,230}

\newcommand{\SOne}{\textcolor{s1color}{\textbf{[S1]}}}
\newcommand{\STwo}{\textcolor{s2color}{\textbf{[S2]}}}
\newcommand{\SThree}{\textcolor{s3color}{\textbf{[S3]}}}

\setlist[itemize]{leftmargin=*, itemsep=1pt, topsep=2pt, parsep=0pt}

\newtcolorbox{keybox}[1]{
  colback=boxbg, colframe=boxborder, boxrule=0.5pt,
  arc=2pt, left=4pt, right=4pt, top=2pt, bottom=2pt,
  title=#1, fonttitle=\bfseries\small,
  coltitle=white, colbacktitle=boxborder
}

\newtcolorbox{formulabox}{
  colback=formulabg, colframe=orange!60!black, boxrule=0.4pt,
  arc=2pt, left=4pt, right=4pt, top=2pt, bottom=2pt
}

\newcommand{\qhead}[2]{%
  \noindent
  \begin{tcolorbox}[colback=boxborder, colframe=boxborder, arc=2pt, left=6pt, right=6pt, top=2pt, bottom=2pt]
    \color{white}\textbf{\large Q#1 -- #2}
  \end{tcolorbox}
}

\newcommand{\huskeord}[1]{%
  \vspace{2pt}
  \begin{tcolorbox}[colback=yellow!15, colframe=orange!70!black, boxrule=0.4pt, arc=2pt, left=4pt, right=4pt, top=2pt, bottom=2pt]
    \footnotesize\textbf{Huskeord:} #1
  \end{tcolorbox}
}

\pagestyle{empty}

\begin{document}

% ==================== Q01 ====================
\qhead{01}{Responsibility centers -- de 5 sentertyper}

\begin{keybox}{Tese}
Centertype = beslutningsrett \SOne\ + måltall \STwo\ + insentiv \SThree. Velg sentertype etter \emph{controllability-prinsippet}: mål dekker det lederen faktisk styrer.
\end{keybox}

\begin{multicols}{2}
\textbf{De 5 typene (stige av kontroll):}
\begin{itemize}
  \item \textbf{Cost center} -- kontrollerer kost (gitt output)
  \item \textbf{Expense center} -- kontrollerer budsjett (output uobs.)
  \item \textbf{Revenue center} -- kontrollerer inntekt
  \item \textbf{Profit center} -- inntekt + kost
  \item \textbf{Investment center} -- profit + kapital
\end{itemize}

\columnbreak

\textbf{Ulemper pr.\ type:}
\begin{itemize}
  \item Cost: kvalitet ofres
  \item Expense: empire building
  \item Revenue: $MR=0 \neq MR=MC$
  \item Profit: transfer pricing-problem
  \item Investment: ROA kortsiktig $\Rightarrow$ EVA
\end{itemize}
\end{multicols}

\begin{formulabox}
\textbf{Controllability:} Måltallet skal dekke \emph{nøyaktig} det lederen kan styre -- ikke mer, ikke mindre.
\end{formulabox}

\textbf{Eks.:} Equinor (Leting = investment, Mongstad = cost). Mærsk PtX = investment center med EVA.

\huskeord{Cost--Expense--Revenue--Profit--Investment $\;\big|\;$ stige $\;\big|\;$ controllability $\;\big|\;$ ROA $\to$ EVA $\;\big|\;$ \STwo}

\vfill\newpage

% ==================== Q02 ====================
\qhead{02}{Transfer Pricing -- costless vs.\ asymmetrisk info}

\begin{keybox}{Tese}
TP $=$ opportunity cost. (a) Costless $\Rightarrow$ first-best. (b) Asymmetrisk $\Rightarrow$ trade-off, ingen metode dominerer.
\end{keybox}

\begin{formulabox}
TP$^{*}$ $=$ MC (ledig kapasitet) $\;\;\big|\;\;$ TP$^{*}$ $=$ MC $+$ fortrengt DB (kapasitetsbinding)
\end{formulabox}

\textbf{5 metoder:}
\begin{itemize}
  \item \textbf{Market-based} -- lik opp.cost, krever likvid marked
  \item \textbf{Marginal cost} -- optimal v/ledig kap., MC overrapporteres
  \item \textbf{Full cost} -- enkel, gir double markup
  \item \textbf{Full cost + markup} -- selger får margin, sterk gaming
  \item \textbf{Negotiated} -- bruker lokal info, forhandlingsstyrke skjevhet
\end{itemize}

\textbf{Asymmetrisk info-problemer:}
\begin{itemize}
  \item Selger overrapporterer MC $\Rightarrow$ for lav intern mengde
  \item Kjøper underdriver WTP $\Rightarrow$ presser pris ned
  \item \emph{Decision management} (riktig mengde) vs.\ \emph{decision control} (riktig insentiv) -- samme tall fungerer ikke til begge
\end{itemize}

\textbf{Eks.:} Equinor (Brent, market-based), Mærsk Ocean/L\&S (cost-plus, friksjon), Mærsk PtX (intet marked $\Rightarrow$ strategisk subsidie).

\huskeord{Opportunity cost $\;\big|\;$ 5 metoder $\;\big|\;$ double markup $\;\big|\;$ mgmt vs.\ control $\;\big|\;$ \STwo}

\vfill\newpage

% ==================== Q03 ====================
\qhead{03}{Costless contracting -- firmaet og de 5 konfliktene}

\begin{keybox}{Tese}
Firma = \emph{nexus of contracts}. Under costless contracting kan alle 5 owner--manager-konflikter kontraktfestes vekk $\Rightarrow$ first-best.
\end{keybox}

\textbf{De 5 konfliktene:}
\begin{enumerate}
  \item \textbf{Shirking} -- innsats under optimum
  \item \textbf{Perks/konsum} -- privat forbruk på firmaets regning
  \item \textbf{Risikoaversjon} -- agent unngår positiv-NPV-risiko
  \item \textbf{Horisont} -- agent ser kortsiktig, eier langsiktig
  \item \textbf{Overinvestering/empire building} -- vekst foran verdi
\end{enumerate}

\begin{formulabox}
Costless contracting: ingen kostnad ved \emph{search, forhandling, drafting, monitoring, enforcement} $\Rightarrow$ kontrakten lukker alle gap.
\end{formulabox}

\textbf{Principal-agent (lineær):} $Q=\alpha e + \mu$, $W = W_0 + \beta Q$. Costless $\Rightarrow$ $\beta$ kan settes uten risiko-trade-off.

\textbf{Spilteori:} monitor vs.\ shirk -- mixed strategy; repetert spill $\Rightarrow$ reputasjon (implicit contracting).

\textbf{Eks.:} Equinor CEO-LTI (tilnærming via vesting). Kontrast Q04: når forutsetningene brytes $\to$ second-best.

\huskeord{Nexus of contracts $\;\big|\;$ 5 konflikter (shirk-perks-risk-horisont-empire) $\;\big|\;$ first-best $\;\big|\;$ \SThree}

\vfill\newpage

% ==================== Q04 ====================
\qhead{04}{Costly contracting + asymmetrisk info -- agency costs}

\begin{keybox}{Tese}
Samme 5 konflikter som Q03, men kontraktskostnader + asymmetrisk info $\Rightarrow$ ufullstendig kontrakt $\Rightarrow$ second-best + residual loss.
\end{keybox}

\textbf{Costly contracting -- 5 kostnadstyper:}
\begin{itemize}
  \item Search, forhandling, drafting, monitoring, enforcement
\end{itemize}

\textbf{To info-problemer:}
\begin{itemize}
  \item \textbf{Adverse selection} (pre-kontrakt, skjult \emph{type}) $\Rightarrow$ signaling, screening, due diligence
  \item \textbf{Moral hazard} (post-kontrakt, skjult \emph{handling}) $\Rightarrow$ monitoring, bonding, insentiv ($\beta > 0$)
\end{itemize}

\begin{formulabox}
\textbf{Agency costs} $=$ Monitoring $+$ Bonding $+$ Residual loss $\;\;$(Jensen \& Meckling 1976)
\end{formulabox}

\textbf{Lineær PA-modell:} optimal $\beta < 1$ balanserer motivasjon mot risikopremie + agency cost.

\textbf{Eks.:} Statoil-Hydro 2007 (adverse selection $\to$ due-diligence). DNB rådgivere (moral hazard $\to$ regulert $\beta$). Mærsk PtX: teknologipartner = AS, EPC = MH.

\huskeord{Adverse selection (type, pre) $\;\big|\;$ Moral hazard (handling, post) $\;\big|\;$ MBR = agency cost $\;\big|\;$ second-best $\;\big|\;$ \SThree}

\vfill\newpage

% ==================== Q05 ====================
\qhead{05}{Organizational architecture -- 3-benet skammel}

\begin{keybox}{Tese}
Arkitektur = \SOne\ + \STwo\ + \SThree\ designet \emph{samtidig}. Endring i ett ben krever justering av de andre, ellers eksploderer agentomkostninger.
\end{keybox}

\textbf{3 ben:}
\begin{itemize}
  \item \SOne\ Decision rights -- hvem bestemmer hva
  \item \STwo\ Performance evaluation -- hva måles
  \item \SThree\ Reward system -- hvordan belønnes
\end{itemize}

\textbf{Begrep a -- Agentomkostninger:} Monitoring + Bonding + Residual loss. Aksepter residual loss (billigere enn perfekt kontroll).

\textbf{Begrep b -- Spesifik vs.\ generell viden:} Hayek $\to$ J\&M. Spesifik viden $\Rightarrow$ desentraliser $\Rightarrow$ krever \STwo\ + \SThree.

\textbf{Begrep c -- Insentiver:} vri egennytte mot prinsipal. Krever målbarhet, kontrollerbarhet, informativeness. Risk-incentive trade-off.

\begin{formulabox}
Komplementer: $\partial^{2} \Pi / \partial S_i \partial S_j > 0$. Misalignment når ett ben endres alene.
\end{formulabox}

\textbf{Eks.:} Equinor Renewables-utskillelse 2020/21 (alle 3 ben endret). Mærsk PtX investment center.

\huskeord{3-benet skammel $\;\big|\;$ MBR $\;\big|\;$ Hayek $\;\big|\;$ komplementer $\;\big|\;$ \SOne+\STwo+\SThree}

\vfill\newpage

% ==================== Q06 ====================
\qhead{06}{Sentralisering vs.\ desentralisering -- optimal grad}

\begin{keybox}{Tese}
Optimal $D^{*}$ der marginal fordel av lokal viden $=$ marginal kostnad av agent- og koordineringsproblemer.
\end{keybox}

\begin{formulabox}
$\mathrm{NB}(D) = BD - AD - CD^{2}$ \\
FOC: $B - A - 2CD = 0 \;\Rightarrow\; D^{*} = (B-A)/(2C)$ \\
SOC: $-2C < 0$ (maks)
\end{formulabox}

\textbf{Komparativ statikk:}
\begin{itemize}
  \item $\partial D^{*}/\partial B > 0$ -- mer lokal viden $\Rightarrow$ desentraliser
  \item $\partial D^{*}/\partial A < 0$ -- mer agentkost $\Rightarrow$ sentraliser
  \item $\partial D^{*}/\partial C < 0$ -- mer koord.kost $\Rightarrow$ sentraliser
\end{itemize}

\textbf{Driver:} spesifik vs.\ generell viden (Hayek). \emph{Co-location of knowledge and decision rights}.

\textbf{Fordeler desentralisering:} lokal viden, hurtighet, toppledertid, motivasjon. \textbf{Ulemper:} agentkost, koordinasjon, tap av stordrift.

\textbf{Eks.:} Equinor E\&P (høy $D^{*}$) vs.\ konsernfinans (lav). IKEA (sentral design, lokal salg). McDonald's franchise. Mærsk PtX: tek desentralt, CAPEX sentralt.

\huskeord{$D^{*}=(B-A)/(2C)$ $\;\big|\;$ Hayek $\;\big|\;$ co-location $\;\big|\;$ \SOne\ flytter \STwo+\SThree}

\vfill\newpage

% ==================== Q07 ====================
\qhead{07}{Decision management vs.\ decision control}

\begin{keybox}{Tese}
Beslutninger har 4 trinn. \emph{Mgmt} $=$ 1+3, \emph{Control} $=$ 2+4. Separer når beslutningstaker $\neq$ residual claimant.
\end{keybox}

\textbf{4 trinn (Fama \& Jensen 1983):}
\begin{enumerate}
  \item \textbf{Initiation} (mgmt)
  \item \textbf{Ratification} (control)
  \item \textbf{Implementation} (mgmt)
  \item \textbf{Monitoring} (control)
\end{enumerate}

\textbf{(a) Allokering:} co-location -- mgmt der spesifik viden er, control der residualkravshaver representeres.

\textbf{(b) Eksternaliteter:} lokale beslutninger med konsernkonsekvenser krever ratification høyere oppe. Internalisering via interne priser/sentral ratifikasjon/KPI.

\textbf{(c) Residual claimants:} driveren.
\begin{itemize}
  \item Sole proprietor: ingen separasjon
  \item Børsnotert: full separasjon
  \item Franchise: lokal residual $\Rightarrow$ delvis desentralisering OK
\end{itemize}

\textbf{Eks.:} Equinor styre (ratification + monitoring). Wells Fargo 2016 (monitoring kollapset $\to$ 3,5 mill.\ falske kontoer). Mærsk PtX: leder initierer, styret ratifiserer CAPEX.

\huskeord{IRIM (Init-Ratify-Impl-Monitor) $\;\big|\;$ Fama \& Jensen $\;\big|\;$ separasjon disiplinerer $\;\big|\;$ \SOne}

\vfill\newpage

% ==================== Q08 ====================
\qhead{08}{Tiltrekke og kompensere ansatte}

\begin{keybox}{Tese}
Sett sammen pakke ($S$, $F$, ikke-monetær) som passerer reservation utility $U_0$ til \emph{lavest kostnad}. Optimum = tangering indifferens $\leftrightarrow$ isokost.
\end{keybox}

\textbf{Mål med kompensasjon:} \emph{tiltrekke, retentere, motivere}.

\begin{formulabox}
$\min c_S S + c_F F \quad \text{s.t.} \quad U(S,F) \geq U_0$ \\
Optimum: $\mathrm{MRS}_{S,F} = c_S / c_F$
\end{formulabox}

\textbf{Kjernebegreper:}
\begin{itemize}
  \item \textbf{Reservation utility $U_0$} -- gulv, neste-beste alternativ, multidim., participation constraint
  \item \textbf{Compensating differentials} -- tillegg for risiko/skift/isolasjon (Smith $\to$ Rosen 1974)
  \item \textbf{Superior assets} -- firmaspesifikke fordeler som senker $c_F$ (skatt, skala, brand, karriere)
  \item \textbf{Salary--fringe-mix} -- konvekse indifferenskurver, rett isokost
\end{itemize}

\textbf{Eks.:} Equinor (skattefordel pensjon), Mærsk shipping-tillegg for sjøtid, Google-perks (kultur som superior asset).

\huskeord{Tiltrekke-retentere-motivere $\;\big|\;$ $U_0$ $\;\big|\;$ MRS = $c_S/c_F$ $\;\big|\;$ comp.\ diff.\ $\;\big|\;$ superior assets $\;\big|\;$ \SThree}

\vfill\newpage

% ==================== Q09 ====================
\qhead{09}{Principal-Agent -- innsats og insentivintensitet}

\begin{keybox}{Tese}
To FOC i samme lineære PA-modell: agentens for $e^{*}$, prinsipalens for $\beta^{*}$.
\end{keybox}

\textbf{Setup:} $Q = \alpha e + \mu$, $W = W_0 + \beta Q$, agent CARA + $c(e)$, prinsipal diversifisert.

\begin{formulabox}
\textbf{Agentens FOC:} $\beta \alpha = c'(e^{*})$ \\
\textbf{Prinsipalens FOC:} $\displaystyle \beta^{*} = \frac{1}{1 + r \sigma^{2} c'' / \alpha^{2}}$
\end{formulabox}

\textbf{(a) Risiko:} $r, \sigma^{2} \uparrow \;\Rightarrow\; \beta^{*} \downarrow$. Risikopremie $= \tfrac{r}{2}\beta^{2}\sigma^{2}$.

\textbf{(b) $W_0$ (fast lønn):} \emph{vrir ikke innsats}. Rolle: IR-constraint + risikopremie-kompensasjon.

\textbf{(c) $\beta$ (variabel andel):} \emph{eneste} innsatsdriver = pay-performance sensitivity.

\textbf{Determinanter for $\beta^{*}$:} $\alpha \uparrow$, $r \downarrow$, $\sigma^{2} \downarrow$, $c'' \downarrow \;\Rightarrow\; \beta^{*} \uparrow$.

\textbf{Eks.:} IB-trader (høy $\beta$, lav $W_0$) vs.\ Equinor-CEO (moderat $\beta$, høy $W_0$). Mærsk PtX: høy $\sigma^{2}$ $\Rightarrow$ lav $\beta$ på EBITDA, moderat på milestones.

\huskeord{$\beta\alpha = c'(e)$ $\;\big|\;$ $\beta^{*} = 1/(1+r\sigma^{2}c''/\alpha^{2})$ $\;\big|\;$ $W_0$ vrir ikke innsats $\;\big|\;$ \STwo+\SThree}

\vfill\newpage

% ==================== Q10 ====================
\qhead{10}{Performance measurement -- 4 formål og 4 problemer}

\begin{keybox}{Tese}
Måling tjener 4 formål samtidig $\Rightarrow$ konflikter $\Rightarrow$ problemene i (a)--(d). Løs med smartere systemdesign, ikke ``bedre måling''.
\end{keybox}

\textbf{4 formål:}
\begin{itemize}
  \item Feedback, Koordinasjon, Belønning/sanksjon, Arkitekturstøtte
\end{itemize}

\textbf{(a) Ratchet effect:} mål fra historie $\Rightarrow$ sandbagging. (Sovjet). Løs: flerårig LTI, ekstern benchmark, lineær pay (Jensen).

\textbf{(b) Gaming + horisont:} Goodhart, multitasking (Holmström \& Milgrom 1991). Løs: BSC, bonusbank, claw-back, subjektiv overstyring.

\textbf{(c) Absolutt vs.\ relativ:} Holmström '79 -- relativ filtrerer felles støy, men inviterer sabotasje. Hybrid (rTSR).

\textbf{(d) Subjektiv vs.\ objektiv:} obj.\ verifiserbar men gaming-utsatt; subj.\ fanger multidim.\ men gir påvirkningskost + bias. Hybrid.

\begin{formulabox}
Goodhart's law: ``When a measure becomes a target, it ceases to be a good measure.''
\end{formulabox}

\textbf{Eks.:} Wells Fargo 2016 (obj.+abs.+høy $\beta$ $\to$ falske kontoer). Microsoft stack ranking $\to$ ``Connect''. Mærsk PtX: milestones + subj.\ + rel.\ peer-benchmark.

\huskeord{Feedback-Koord-Belønning-Arkitektur $\;\big|\;$ Ratchet $\;\big|\;$ Gaming $\;\big|\;$ rTSR $\;\big|\;$ \STwo}

\vfill\newpage

% ==================== Q11 ====================
\qhead{11}{Vertikal integrasjon vs.\ outsourcing}

\begin{keybox}{Tese}
Firmagrenser bestemmes av transaksjonskost vs.\ byråkratikost. VI lønner seg ved spesifikke aktiva, usikkerhet, markedsmakt.
\end{keybox}

\textbf{Teoretisk fundament:} Coase, Williamson, Grossman-Hart-Moore (eierskap = residuell kontroll).

\begin{formulabox}
\textbf{Economic profit:} $\pi_{\text{econ}} = R - C_{\text{exp}} - C_{\text{opp}}$ \\
(Skiller seg fra accounting profit ved alternativkostnad.)
\end{formulabox}

\textbf{Fordeler outsourcing:} markedsdisiplin, skala, fleksibilitet, sterkere insentiver.

\textbf{Fordeler VI:} hold-up eliminert, kontroll, koordinasjon, IP-beskyttelse, market power.

\textbf{Market power-modell:} $U$ monopol $\to$ integrerer $D_1$ $\to$ foreclosure av $D_2$.

\textbf{Chicago-kritikk:} VI tilfører ikke ny markedsmakt hvis $U$ allerede ekstraherer via inputpris (\emph{one monopoly rent theorem}).

\textbf{Eks.:} De Beers (foreclosure). Mærsk/APM Terminals (effektivitet + delvis foreclosure). Mærsk PtX: make-or-buy elektrolyse.

\huskeord{Coase-Williamson-GHM $\;\big|\;$ Acc.\ vs.\ econ.\ profit $\;\big|\;$ foreclosure $\;\big|\;$ Chicago $\;\big|\;$ \SOne+\STwo}

\vfill\newpage

% ==================== Q12 ====================
\qhead{12}{Outsourcing -- hold-up, free-rider, double markup}

\begin{keybox}{Tese}
Tre kontraktproblemer forklarer når outsourcing feiler og VI lønner seg.
\end{keybox}

\textbf{1. Hold-up:} spesifikke aktiva $\Rightarrow$ quasi-rente approprieres $\Rightarrow$ \emph{underinvestering}.
\begin{itemize}
  \item 5 typer asset specificity: site, physical, human, dedicated, brand
  \item Løs: VI, langsiktig kontrakt, mutual hostage
\end{itemize}

\textbf{2. Free-rider:} distributør underinvesterer i service (\emph{showrooming}).
\begin{itemize}
  \item Løs: franchise, eksklusive territorier, resale price maintenance (RPM)
\end{itemize}

\textbf{3. Double markup:} hvert ledd legger på margin $\Rightarrow$ pris for høy, kvantum for lavt.

\begin{formulabox}
Eks.: $c=20$, $P_w=60$ (oppstrøms), $P=80$, $Q=20$ \\
Integrert: $P=60$, $Q=40$ \\
Løs: VI, two-part tariff (fee + per-unit)
\end{formulabox}

\textbf{Eks.:} IKEA rammeavtaler (hold-up), McDonald's franchise (free-rider), Mærsk PtX (alle tre problemene representert).

\huskeord{Hold-up (spes.aktiva) $\;\big|\;$ Free-rider (service) $\;\big|\;$ Double markup (margin-stacking) $\;\big|\;$ \SOne+\SThree}

\vfill\newpage

% ==================== Q13 ====================
\qhead{13}{Regulering og det politiske støttemarkedet}

\begin{keybox}{Tese}
Regulering hensiktsmessig ved 4 markedssvikter, men kun hvis reguleringsgevinst $>$ reguleringskostnad. $P^{*}$ er politisk kompromiss, ikke velferdsmaks.
\end{keybox}

\textbf{4 markedssvikter:}
\begin{enumerate}
  \item Naturlig monopol ($AC$ fallende)
  \item Eksternaliteter (Pigou)
  \item Offentlige goder (free-rider)
  \item Informasjonsasymmetri (lemons)
\end{enumerate}

\textbf{Politisk støttemodell (Stigler/Peltzman):}
\begin{formulabox}
Maks $S(P) = f(\pi(P), CS(P))$ \\
$\pi(P)$ invers-U (topp ved $P^M$); $CS(P)$ fallende \\
$P^{*}$ der bransjens marginale pull opp $=$ konsumentens pull ned
\end{formulabox}

\textbf{Komparativ statikk:} sterkere lobby $\Rightarrow$ $P^{*} \to P^{M}$; sterkere konsument $\Rightarrow$ $P^{*} \to P^{MC}$.

\textbf{Kritikk:} government failure, capture (Stigler), Averch-Johnson overinvestering, Olsons kollektive handling.

\textbf{Eks.:} Norsk farmasiregulering (monopol + lemons). EU ETS (eksternalitet). Mærsk PtX (subsidier + IMO Pigou).

\huskeord{4 svikter $\;\big|\;$ Stigler/Peltzman $\;\big|\;$ capture $\;\big|\;$ AJ $\;\big|\;$ $P^{*}$ = kompromiss $\;\big|\;$ \SOne}

\vfill\newpage

% ==================== Q14 ====================
\qhead{14}{Informativeness-prinsippet}

\begin{keybox}{Tese}
Holmström (1979): inkluder signal $y$ i kontrakten \emph{bare} hvis det tilfører informasjon om $e$ utover $Q$. Sufficient statistic.
\end{keybox}

\begin{formulabox}
$y$ informativt $\Leftrightarrow$ posterior for $e$ endres gitt $(Q, y)$ vs.\ gitt $Q$ alene. \\
$\beta^{*} = \dfrac{1}{1 + r \sigma^{2} c'' / \alpha^{2}}$
\end{formulabox}

\textbf{(a) Kontrollerbarhet:} spesialtilfelle av informativeness; ukontrollerbare mål kan filtrere felles støy.

\textbf{(b) Risikodeling:} bedre signaler $\Rightarrow$ Pareto-bedre risk-sharing for gitt innsats.

\textbf{(c) Presisjon:} lav $\sigma^{2}$ $\Rightarrow$ høyere $\beta^{*}$ -- mer pay-for-performance når støy er lav.

\textbf{(d) Risikopremie:} $RP = \tfrac{r}{2}\beta^{2}\sigma^{2}$; lavere effektiv $\sigma^{2}$ $\Rightarrow$ lavere RP.

\textbf{Hovedeksempel:} \emph{Relativ TSR (rTSR)} -- peer-signal filtrerer markedssjokk.

\textbf{Kritikk:} multitasking (Holmström \& Milgrom 1991), kognitiv overbelastning, informasjonskostnad.

\textbf{Eks.:} Equinor CEO med rTSR. Sykehuslege (upresise mål $\to$ subjektiv evaluering). Mærsk PtX milestones + peer-benchmark.

\huskeord{Holmström '79 $\;\big|\;$ sufficient statistic $\;\big|\;$ rTSR $\;\big|\;$ presisjon $\to$ $\beta\uparrow$ $\;\big|\;$ \STwo+\SThree}

\vfill\newpage

% ==================== Q15 ====================
\qhead{15}{Nyttemodeller -- integritet og risiko}

\begin{keybox}{Tese}
To nyttemodeller: (1) integritet vs.\ inntekt -- kompensasjon påvirker etisk atferd; (2) risiko vs.\ forventet verdi -- risikoaversjon bestemmer risikopremie.
\end{keybox}

\textbf{Modell 1 -- Integritet vs.\ inntekt:}
\begin{itemize}
  \item Akser: integritet (x), inntekt (y)
  \item Konvekse indifferenskurver; budsjettlinje $=$ kompensasjonsplan
  \item Aggressivt insentiv $\Rightarrow$ rotasjon av budsjett $\Rightarrow$ A $\to$ B (mindre integritet)
  \item Innsikt: \emph{system, ikke person}. $\beta$ driver rotasjonen.
\end{itemize}

\textbf{Modell 2 -- Risiko vs.\ E[W]:}
\begin{formulabox}
Risikoavers indifferens: $E[W] = \bar{U} + \tfrac{r}{2}\sigma^{2}$ \\
Certainty equivalent: $CE = E[W] - \tfrac{r}{2}\sigma^{2}$ \\
Risikopremie: $RP = \tfrac{r}{2}\beta^{2}\sigma^{2}$
\end{formulabox}

\begin{itemize}
  \item Risikoavers: stigende, konveks indifferenskurve
  \item Risikonøytral: flat indifferenskurve
\end{itemize}

\textbf{Kobling:} $\beta$ er felles designvariabel. Høy $\beta$ $\Rightarrow$ (1) mer fristelse + (2) mer risiko.

\textbf{Eks.:} Wells Fargo 2016 (modell 1 -- falske kontoer). Equinor aksjeprogram (modell 2). Mærsk PtX: prognoseinflasjon + lederrekruttering med høy $\sigma$.

\huskeord{Modell 1: integritet--inntekt $\;\big|\;$ Modell 2: $E[W]$--$\sigma$ $\;\big|\;$ $RP = \tfrac{r}{2}\beta^{2}\sigma^{2}$ $\;\big|\;$ system $\neq$ person $\;\big|\;$ \SThree}

\end{document}
